\section{Base change}
\label{section-base-change}

\noindent
We formally introduce base change in algebra as follows.

\begin{definition}
\label{definition-base-change}
Let $\varphi : R \to S$ be a ring map. Let $M$ be an $S$-module.
Let $R \to R'$ be any ring map. The {\it base change} of $\varphi$
by $R \to R'$ is the ring map $R' \to S \otimes_R R'$. In this situation
we often write $S' = S \otimes_R R'$.
The {\it base change} of the $S$-module $M$ is the $S'$-module
$M \otimes_R R'$.
\end{definition}

\noindent
If $S = R[x_i]/(f_j)$ for some collection of variables $x_i$, $i \in I$
and some collection of polynomials $f_j \in R[x_i]$, $j \in J$, then
$S \otimes_R R' = R'[x_i]/(f'_j)$, where $f'_j \in R'[x_i]$ is the image
of $f_j$ under the map $R[x_i] \to R'[x_i]$ induced by $R \to R'$.
This simple remark is the key to understanding base change.

\begin{lemma}
\label{lemma-base-change-finiteness}
Let $R \to S$ be a ring map. Let $M$ be an $S$-module.
Let $R \to R'$ be a ring map and let $S' = S \otimes_R R'$ and
$M' = M \otimes_R R'$ be the base changes.
\begin{enumerate}
\item If $M$ is a finite $S$-module, then the base change
$M'$ is a finite $S'$-module.
\item If $M$ is an $S$-module of finite presentation, then
the base change $M'$ is an $S'$-module of finite presentation.
\item If $R \to S$ is of finite type, then the base change
$R' \to S'$ is of finite type.
\item If $R \to S$ is of finite presentation, then
the base change $R' \to S'$ is of finite presentation.
\end{enumerate}
\end{lemma}

\begin{proof}
Proof of (1). Take a surjective, $S$-linear map
$S^{\oplus n} \to M \to 0$.
By Lemma \ref{lemma-flip-tensor-product} and \ref{lemma-tensor-product-exact}
the result after tensoring with $R^\prime$ is a surjection
${S^\prime}^{\oplus n} \to M^\prime \rightarrow 0$,
so $M^\prime$ is a finitely generated $S^\prime$-module.
Proof of (2). Take a presentation
$S^{\oplus m} \to S^{\oplus n} \to M \to 0$.
By Lemma \ref{lemma-flip-tensor-product} and \ref{lemma-tensor-product-exact}
the result after tensoring with $R^\prime$ gives a finite presentation
${S^\prime}^{\oplus m} \to {S^\prime}^{\oplus n} \to M^\prime \to 0$, of
the $S^\prime$-module $M^\prime$. Proof of (3). This follows by the remark
preceding the lemma as we can take $I$ to be finite by assumption.
Proof of (4). This follows by the remark preceding the lemma
as we can take $I$ and $J$ to be finite by assumption.
\end{proof}

\noindent
Let $\varphi : R \to S$ be a ring map. Given an $S$-module $N$ we
obtain an $R$-module $N_R$ by the rule $r \cdot n = \varphi(r)n$.
This is sometimes called the {\it restriction} of $N$ to $R$.

\begin{lemma}
\label{lemma-adjoint-tensor-restrict}
Let $R \to S$ be a ring map. The functors
$\text{Mod}_S \to \text{Mod}_R$, $N \mapsto N_R$ (restriction)
and $\text{Mod}_R \to \text{Mod}_S$, $M \mapsto M \otimes_R S$
(base change) are adjoint functors. In a formula
$$
\Hom_R(M, N_R) = \Hom_S(M \otimes_R S, N)
$$
\end{lemma}

\begin{proof}
If $\alpha : M \to N_R$ is an $R$-module map, then we define
$\alpha' : M \otimes_R S \to N$ by the rule
$\alpha'(m \otimes s) = s\alpha(m)$. If $\beta : M \otimes_R S \to N$
is an $S$-module map, we define $\beta' : M \to N_R$ by the rule
$\beta'(m) = \beta(m \otimes 1)$.
We omit the verification that these constructions are mutually inverse.
\end{proof}

\noindent
The lemma above tells us that restriction has a left adjoint, namely
base change. It also has a right adjoint.

\begin{lemma}
\label{lemma-adjoint-hom-restrict}
Let $R \to S$ be a ring map. The functors
$\text{Mod}_S \to \text{Mod}_R$, $N \mapsto N_R$ (restriction)
and $\text{Mod}_R \to \text{Mod}_S$, $M \mapsto \Hom_R(S, M)$
are adjoint functors. In a formula
$$
\Hom_R(N_R, M) = \Hom_S(N, \Hom_R(S, M))
$$
\end{lemma}

\begin{proof}
If $\alpha : N_R \to M$ is an $R$-module map, then we define
$\alpha' : N \to \Hom_R(S, M)$ by the rule
$\alpha'(n) = (s \mapsto \alpha(sn))$. If $\beta : N \to \Hom_R(S, M)$
is an $S$-module map, we define $\beta' : N_R \to M$ by the rule
$\beta'(n) = \beta(n)(1)$.
We omit the verification that these constructions are mutually inverse.
\end{proof}

\begin{lemma}
\label{lemma-hom-from-tensor-product-variant}
Let $R \to S$ be a ring map. Given $S$-modules $M, N$ and an $R$-module $P$
we have
$$
\Hom_R(M \otimes_S N, P) = \Hom_S(M, \Hom_R(N, P))
$$
\end{lemma}

\begin{proof}
This can be proved directly, but it is also a consequence of
Lemmas \ref{lemma-adjoint-hom-restrict} and \ref{lemma-hom-from-tensor-product}.
Namely, we have
\begin{align*}
\Hom_R(M \otimes_S N, P)
& =
\Hom_S(M \otimes_S N, \Hom_R(S, P)) \\
& =
\Hom_S(M, \Hom_S(N, \Hom_R(S, P))) \\
& =
\Hom_S(M, \Hom_R(N, P))
\end{align*}
as desired.
\end{proof}








